\documentclass[10pt,aspectratio=169]{beamer}
% Preámbulo modular para presentaciones Beamer (Modelación Estadística)
% Diseño y paleta institucional del Tecnológico de Monterrey

\documentclass[aspectratio=169,xcolor={dvipsnames,table}]{beamer}

\usepackage[utf8]{inputenc}
\usepackage[spanish,mexico]{babel}
\usepackage{amsmath,amssymb,amsfonts,mathtools}
\usepackage{graphicx}
\usepackage{listings}
\usepackage{booktabs}
\usepackage{tikz}
\usetikzlibrary{matrix,arrows,decorations.pathmorphing,shapes.geometric,positioning}

% Paleta Institucional Tecnológico de Monterrey
\definecolor{TecAzul}{HTML}{0039A6}       % Azul TEC: color primario institucional
\definecolor{TecAzulOscuro}{HTML}{1A2E51} % Azul marino oscuro
\definecolor{TecRojo}{HTML}{EC2661}       % Rojo de acento (paleta miTec)
\definecolor{TecGris}{HTML}{646464}       % Gris neutro para comentarios y subtítulos
\definecolor{TecFondoBloque}{HTML}{F4F6F9}% Fondo suave para bloques

% Configuración de Tema Beamer
\usetheme{Boadilla}
\usecolortheme[named=TecAzulOscuro]{structure}
\setbeamercolor{alerted text}{fg=TecRojo}
\setbeamercolor{palette primary}{bg=TecAzulOscuro,fg=white}
\setbeamercolor{palette secondary}{bg=TecAzul,fg=white}
\setbeamercolor{palette tertiary}{bg=TecAzulOscuro!80!black,fg=white}
\setbeamercolor{title}{fg=TecAzulOscuro,bg=TecFondoBloque}
\setbeamercolor{frametitle}{fg=TecAzulOscuro,bg=TecFondoBloque}

% Bloques personalizados elegantes
\setbeamercolor{block title}{bg=TecAzulOscuro,fg=white}
\setbeamercolor{block body}{bg=TecFondoBloque,fg=black}
\setbeamercolor{block title alerted}{bg=TecRojo,fg=white}
\setbeamercolor{block body alerted}{bg=TecRojo!8,fg=black}
\setbeamercolor{block title example}{bg=TecAzul,fg=white}
\setbeamercolor{block body example}{bg=TecAzul!8,fg=black}

% Redondear bordes y añadir sombra sutil
\setbeamertemplate{blocks}[rounded][shadow=true]
\setbeamertemplate{navigation symbols}{}

% Pie de página limpio con número de diapositiva
\setbeamertemplate{footline}{
  \leavevmode%
  \hbox{%
  \begin{beamercolorbox}[wd=.7\paperwidth,ht=2.25ex,dp=1ex,left]{author in head/foot}%
    \hspace*{2ex}\insertshortauthor~~(\insertshortinstitute) \hfill \insertshorttitle
  \end{beamercolorbox}%
  \begin{beamercolorbox}[wd=.3\paperwidth,ht=2.25ex,dp=1ex,right]{date in head/foot}%
    \insertshortdate{}\hspace*{2em}
    \insertframenumber{} / \inserttotalframenumber\hspace*{2ex}
  \end{beamercolorbox}}%
  \vskip0pt%
}

% Configuración de Listings de Python para visualización pedagógica de código
\lstset{
    language=Python,
    basicstyle=\ttfamily\footnotesize,
    keywordstyle=\color{TecAzulOscuro}\bfseries,
    stringstyle=\color{TecRojo},
    commentstyle=\color{TecGris}\itshape,
    showstringspaces=false,
    breaklines=true,
    frame=lines,
    rulecolor=\color{TecAzulOscuro!30},
    backgroundcolor=\color{TecFondoBloque},
    aboveskip=0.5em,
    belowskip=0.5em,
    literate={á}{{\'a}}1 {é}{{\'e}}1 {í}{{\'i}}1 {ó}{{\'o}}1 {ú}{{\'u}}1
             {Á}{{\'A}}1 {É}{{\'E}}1 {Í}{{\'I}}1 {Ó}{{\'O}}1 {Ú}{{\'U}}1
             {ñ}{{\~n}}1 {Ñ}{{\~N}}1 {¿}{{?`}}1 {¡}{{!`}}1
}

% Metadatos institucionales por defecto
\title[Modelación Estadística]{Modelación Estadística para Ciencia de Datos e Ingeniería}
\author[J. Castillo Colmenares]{Juliho David Castillo Colmenares}
\institute[Tec de Monterrey]{Tecnológico de Monterrey \\ Escuela de Ingeniería y Ciencias}
\date{\today}

% Comandos y macros estadísticas compartidas para las presentaciones Beamer (Metropolis)

% Conjuntos numéricos básicos
\newcommand{\R}{\mathbb{R}}
\newcommand{\N}{\mathbb{N}}
\newcommand{\Q}{\mathbb{Q}}
\newcommand{\Z}{\mathbb{Z}}
\newcommand{\set}[1]{\left\{ #1 \right\}}
\newcommand{\abs}[1]{\left|#1\right|}
\newcommand{\norm}[1]{\left\| #1 \right\|}

% Comandos estadísticos y probabilísticos del curso
\newcommand{\Var}{\operatorname{Var}}
\newcommand{\cov}{\operatorname{Cov}}
\newcommand{\corr}{\rho}
\newcommand{\comb}[2]{\begin{pmatrix} #1 \\ #2 \end{pmatrix}}
\newcommand{\s}{\sigma}
\newcommand{\particion}{\mathfrak{P}}
\newcommand{\card}[1]{\#\left( #1 \right)}
\newcommand{\seq}[3]{\set{#1}_{#2}^{#3}}
\newcommand{\E}{\mathbb{E}}
\newcommand{\Prob}{\mathbb{P}}

% Entornos pedagógicos y cajas en español académico natural
\newenvironment{discusion}{%
  \begin{alertblock}{Pregunta de Discusión}%
}{%
  \end{alertblock}%
}

\newenvironment{reflexion}{%
  \begin{alertblock}{Reflexión para el Aula}%
}{%
  \end{alertblock}%
}

\newenvironment{paradoja}{%
  \begin{alertblock}{Paradoja de Exploración}%
}{%
  \end{alertblock}%
}

\newenvironment{motivacion}{%
  \begin{exampleblock}{Motivación: Ciencia de Datos en la Práctica}%
}{%
  \end{exampleblock}%
}

\newenvironment{retoclase}{%
  \begin{block}{Reto Rápido en Equipo}%
}{%
  \end{block}%
}


\title[Estimación por intervalo]{Sección 6.2: Estimación por intervalo}
\subtitle{Intervalos de confianza, niveles y valores $p$}
\author[J. Castillo Colmenares]{Juliho Castillo Colmenares}
\institute{Tecnológico de Monterrey}
\date{}

\begin{document}
\begin{frame}[plain]\vspace{-0.8cm}\titlepage\end{frame}

\begin{frame}{Hoja de ruta --- capítulo 6}
  \footnotesize
  \begin{itemize}\itemsep=0.08cm
    \item \textbf{06.01} Estimación puntual
    \item \textbf{06.02} Estimación por intervalo \emph{(hoy)}
    \item \textbf{06.04--06.07} Medias, proporciones, varianzas y tamaño muestral
  \end{itemize}
  \begin{block}{Objetivo didáctico}
    Construir intervalos para una media, distinguir $Z$ de $t$ y relacionar
    confianza, significación y valores $p$.
  \end{block}
\end{frame}

\begin{frame}{Motivación: un rango plausible}
  \footnotesize
  Un estimador puntual no muestra toda la incertidumbre. Un intervalo de
  confianza proporciona un rango de valores compatibles con los datos.
  \begin{alertblock}{Pregunta guía}
    ¿Qué tan ancho debe ser el rango para alcanzar la confianza elegida?
  \end{alertblock}
\end{frame}

\begin{frame}{Confianza y valor $p$}
  \footnotesize
  Para un nivel de confianza $p$, la nota define la significación como
  \[
    \beta=1-p.
  \]
  Un valor $p$ de un estadístico se expresa como
  \[
    p\text{-valor}=P(X>Z)\quad\text{o}\quad P(X>t).
  \]
\end{frame}

\begin{frame}{Criterio y tipos de cola}
  \footnotesize
  \begin{itemize}\itemsep=0.12cm
    \item Si $p\text{-valor}>\beta$, no se rechaza la hipótesis nula.
    \item Si $p\text{-valor}<\beta$, se favorece la alternativa.
  \end{itemize}
  La simetría de la normal permite pruebas de cola izquierda, derecha o dos
  colas.
\end{frame}

\begin{frame}{Intervalo para $\mu$ con $\sigma$ conocida}
  \footnotesize
  \[
    \bar X-Z_{\alpha/2}\frac{\sigma}{\sqrt n}
    \le\mu\le
    \bar X+Z_{\alpha/2}\frac{\sigma}{\sqrt n}.
  \]
  El cuantil pertenece a la normal estándar.
\end{frame}

\begin{frame}{Intervalo para $\mu$ con $\sigma$ desconocida}
  \footnotesize
  \[
    \bar X-t_{\alpha/2,n-1}\frac{S}{\sqrt n}
    \le\mu\le
    \bar X+t_{\alpha/2,n-1}\frac{S}{\sqrt n}.
  \]
  Sustituir $\sigma$ por $S$ cambia la distribución de referencia a $t$.
\end{frame}

\begin{frame}{Estructura común}
  \footnotesize
  \[
    \text{estimador puntual}\ \pm\ (\text{valor crítico})
    (\text{error estándar}).
  \]
  Esta forma se repite para medias, proporciones y varianzas.
\end{frame}

\begin{frame}{Puente computacional 1/4: margen de error}
  \footnotesize
  \[
    E=\text{valor crítico}\times SE.
  \]
  Aumentar el nivel de confianza aumenta el valor crítico y ensancha el
  intervalo.
\end{frame}

\begin{frame}{Puente computacional 2/4: Z frente a t}
  \footnotesize
  \begin{tabular}{lll}
    \hline
    Caso & Error estándar & Referencia \\
    \hline
    $\sigma$ conocida & $\sigma/\sqrt n$ & $Z$ \\
    $\sigma$ desconocida & $S/\sqrt n$ & $t_{n-1}$ \\
    \hline
  \end{tabular}
\end{frame}

\begin{frame}{Puente computacional 3/4: colas}
  \footnotesize
  Una prueba bilateral reparte $\alpha$ entre las dos colas. Una prueba de una
  cola concentra toda la región crítica en un extremo.
  \begin{block}{Consecuencia}
    El tipo de cola debe fijarse antes de interpretar el valor $p$.
  \end{block}
\end{frame}

\begin{frame}{Puente computacional 4/4: interpretación}
  \footnotesize
  Un IC del $95\%$ se construye con $\alpha=0.05$ y cuantiles
  $Z_{0.025}$ o $t_{0.025,n-1}$.
  \begin{alertblock}{Cuidado}
    El intervalo cuantifica el procedimiento repetido; no afirma que un
    parámetro fijo sea aleatorio después de observar los datos.
  \end{alertblock}
\end{frame}

\begin{frame}{Ejemplo resuelto 1/4 --- botellas y $Z$}
  \footnotesize
  Una muestra de $n=16$ botellas tiene $\bar X=500.4$ ml y se conoce
  $\sigma=2$ ml. Construir el IC del $95\%$ para $\mu$.
\end{frame}

\begin{frame}{Ejemplo resuelto 1/4 --- solución}
  \footnotesize
  \[
    SE=2/\sqrt{16}=0.5,\qquad
    IC=500.4\pm1.96(0.5)=[499.42,501.38]\text{ ml}.
  \]
  Se emplea la normal porque $\sigma$ es conocida.
\end{frame}

\begin{frame}{Ejemplo resuelto 2/4 --- botellas y $t$}
  \footnotesize
  Con los mismos datos, suponga $\sigma$ desconocida y $S=2.3$ ml. Use
  $t_{15,0.025}\approx2.131$ para el IC del $95\%$.
\end{frame}

\begin{frame}{Ejemplo resuelto 2/4 --- solución}
  \footnotesize
  \[
    SE=2.3/\sqrt{16}=0.575,\qquad
    IC=500.4\pm2.131(0.575)=[499.18,501.63]\text{ ml}.
  \]
  El intervalo $t$ es más ancho por la incertidumbre adicional.
\end{frame}

\begin{frame}{Aplicación de lectura 3/4 --- cola derecha}
  \footnotesize
  Para una hipótesis de aumento, la región crítica se ubica en la cola derecha.
  \begin{block}{Planteamiento}
    Comparar el valor $p$ con $\beta$ usando la probabilidad a la derecha del
    estadístico observado.
  \end{block}
\end{frame}

\begin{frame}{Aplicación de lectura 3/4 --- decisión}
  \footnotesize
  Si el valor $p$ es menor que $\beta$, la evidencia cae en la región crítica y
  se rechaza $H_0$; de lo contrario, no se rechaza.
\end{frame}

\begin{frame}{Aplicación de lectura 4/4 --- dos colas}
  \footnotesize
  Para una diferencia en cualquier dirección, la prueba bilateral distribuye la
  significación entre ambas colas.
  \begin{block}{Planteamiento}
    Interpretar el intervalo como el conjunto de valores no rechazados por la
    prueba bilateral correspondiente.
  \end{block}
\end{frame}

\begin{frame}{Aplicación de lectura 4/4 --- conexión}
  \footnotesize
  Un valor de $\mu_0$ fuera del IC bilateral produce un valor $p$ menor que
  $\alpha$; si está dentro, el procedimiento no lo rechaza.
  \begin{exampleblock}{Resultado}
    IC y prueba bilateral son dos lecturas del mismo margen de error.
  \end{exampleblock}
\end{frame}

\begin{frame}{Síntesis de la sección}
  \footnotesize
  \begin{itemize}\itemsep=0.12cm
    \item El nivel de confianza fija $\alpha$ y el valor crítico.
    \item $Z$ se usa con $\sigma$ conocida; $t$ con $S$.
    \item El tipo de cola define la región crítica.
    \item IC y valor $p$ conectan estimación y decisión.
  \end{itemize}
\end{frame}

\begin{frame}{Transición curricular}
  \footnotesize
  \begin{block}{Lo que queda establecido}
    Un intervalo combina estimador, distribución de referencia y error estándar.
  \end{block}
  \begin{alertblock}{Siguiente sección}
    \textbf{06.04 Errores estándar}: comparar las fórmulas de precisión para
    medias, diferencias y proporciones.
  \end{alertblock}
\end{frame}
\end{document}
